\documentclass[11pt,a4paper]{article}
\usepackage[utf8]{inputenc}
\usepackage[T1]{fontenc}
\IfFileExists{polish.ldf}{\usepackage[polish]{babel}}{}
\IfFileExists{lmodern.sty}{\usepackage{lmodern}}{}
\usepackage[margin=26mm]{geometry}
\usepackage{amsmath,booktabs}
\usepackage[authoryear,round]{natbib}
\usepackage[hidelinks]{hyperref}
\title{Kandydat: walidacja techniczna urządzeń EEG w EEG Neurolyze}
\author{Mateusz Pomianek}
\date{MSC-EEG-P02 | Protokół techniczny 1.0.0 | 27 listopada 2026}
\begin{document}
\maketitle
\begin{abstract}
Przedstawiono prospektywny protokół kwalifikacji ścieżki pomiarowej Muse S Athena i EEG Neurolyze. Ocenie podlegają kompletność akwizycji, pochodzenie znaczników czasu, mapowanie kanałów i zgodność sygnału z niezależnym odniesieniem. Dowody uporządkowano według jednostki, konfiguracji i poziomu walidacji. Siedem testów oprogramowania wykonano bez urządzenia. Zaplanowana macierz fizyczna pozostaje niezmierzona.
\end{abstract}

\section{Wprowadzenie}
Reprodukowalność EEG zależy od udokumentowania urządzenia, rejestracji, metadanych oraz reguł przetwarzania \citep{pernet2020}. EEG-BIDS opisuje sposób zapisu pomiaru i jego metadanych \citep{pernet2019}. Poprawna struktura zbioru oraz jakość pomiaru odpowiadają różnym pytaniom. MSC-EEG-P01 ocenia zgodność numeryczną oprogramowania. W P02 jednostką kwalifikacji jest konfiguracja obejmująca egzemplarz urządzenia, firmware, preset, sterownik, backend i tryb opóźnienia. Wnioski odnoszą się do faktycznie przebadanych konfiguracji.

\section{Materiał i metody}
Punktem wyjścia jest protokół M03 z repozytorium Neurolyze. Dla każdej obsługiwanej konfiguracji planuje się po trzy próby 60 s na ścieżkę. Przechwycenie surowych powiadomień i akwizycję BrainFlow przeprowadza się w osobnych blokach, z jednym właścicielem połączenia w danej chwili. Kolejność bloków należy równoważyć. Rejestrowane są fazy połączenia, liczba wywołań odbiorczych, liczba poprawnie rozpoznanych ramek, kompletnych próbek, odstępstwa oraz powód przerwania. Nieudaną próbę ujmuje się w mianowniku i zachowuje jej częściowy dziennik.

Mapowanie kanałów ocenia się na stanowisku z niezależną identyfikacją wejść i informacją o referencji. Ocena czasu wymaga niezależnego znacznika lub wzorca, który można powiązać z pomiarem, a ocena amplitudy wymaga wzorca o udokumentowanej niepewności. Badania stanowiskowe są oddzielone od przyszłych prób z udziałem osób. Wartość progów dla czasu, amplitudy i zakresu częstotliwości ustala się według zastosowania przed pozyskaniem danych porównawczych.

\section{Plan analizy}
Pierwszorzędowy wskaźnik: odsetek pełnych sesji technicznych spośród wszystkich rozpoczętych prób dla każdej konfiguracji i ścieżki osobno. Pełna sesja wymaga zakończenia akwizycji, poprawnego archiwum i udokumentowanego pokrycia przedziału na właściwej osi czasu. Dla przebiegów demonstracyjnych raportuje się liczbę sukcesów oraz dokładny przedział dwumianowy. Trzy sukcesy w trzech próbach mają dwustronną dolną granicę 95\% równą około 0,292 i dają ograniczoną informację o niezawodności. Przedział ma interpretację przy założeniu niezależności i stałego prawdopodobieństwa sukcesu.

Dla pomiaru zsynchronizowanego estymuje się relację $t_\mathrm{obs}=a+b t_\mathrm{ref}+\varepsilon$, offset $a$, dryf zegara $(b-1)10^6$ w ppm i rozrzut reszt. Przesunięcie po stronie oprogramowania i odchylenie zegara wymagają odpowiednich punktów odniesienia. Osobno raportuje się percentyle opóźnienia zdarzeń i najdłuższe przerwy odbioru. Następstwo liczników analizuje się w obrębie tej samej charakterystyki i sesji, z obsługą zawijania. Zamrożenie estymandu i progów porównawczych nastąpi po wyborze właściwego scenariusza i wzorca, przed odczytem wyników fizycznych.

\section{Wyniki}
W przypiętym kodzie Neurolyze wykonano siedem testów dotyczących kwalifikacji wejść, powtórzonych znaczników czasu i ciągłości indeksu. Wszystkie przeszły. Testy wykrywają między innymi podanie syntetycznego archiwum jako pomiaru, zachowanie zdublowanych znaczników z danych oznaczonych jako mierzone oraz rozróżnienie wadliwej ramki od przerwy w indeksie. Karta kwalifikacji M03 wskazuje brak wykonanych pomiarów dla pełnej macierzy. W tym etapie nie uruchamiano urządzenia i nie oszacowano parametrów toru fizycznego.

\section{Dyskusja i ograniczenia}
Przerwa powiadomień hosta może wynikać z kolejkowania, opóźnienia wywołania albo zachowania transportu. Powtórzone znaczniki mogą być skutkiem działania źródła programowego. Interpretacja ubytku próbek wymaga ustalenia semantyki liczników i mianownika. Różnica liczby rekordów z oddzielnych sesji dwóch backendów opisuje odrębne przebiegi. Kwalifikacja całej ścieżki wymaga dodatkowo sprawdzenia kanałów, powtarzalności i niezależnego pomiaru czasu. Poprawność pliku EEG-BIDS wymaga uruchomienia walidatora dla konkretnego zbioru.
\bibliographystyle{apalike}
\bibliography{references}
\end{document}
